
% 本章介绍端到端神经算子 DeepRTE。该方法以边界条件、截面系数（$\mu_s, \mu_t$）以及散射核 $p(\bOmega,\bOmega^*)$ 为输入，直接学习其到稳态解场的映射关系。
本章介绍端到端神经算子 DeepRTE（Deep Neural Operator for Radiative Transfer Equation）。该方法以边界条件、截面系数以及散射核为输入，直接学习其到稳态解场的映射关系。

本章首先从辐射输运方程（radiative transport equation，RTE）解的内在结构出发，对输运过程进行可解释的数学分解，并将原问题改写为一个等价的不动点形式，使得可以通过迭代逐步逼近真实解。进一步地，为保证所构造算子具有良好的可推广性（尤其是对不同入射边界条件的适应性），我们引入格林函数表述，在统一的理论框架下同时刻画边界信息与介质响应，从而使该方法不局限于某一类特定的入射边界函数。上述推导均由相应的定理与引理支撑，以确保分解与等价变换的严谨性。

由于分解后得到的部分算子通常难以获得解析形式，本章随后给出 DeepRTE 的网络结构设计：使用神经网络对其中的关键算子（例如与衰减/传播相关的映射）进行参数化近似，并将其嵌入不动点迭代流程中，从而形成用于求解辐射输运方程的残差神经算子架构。

